% Generated by blueprint/scripts/audit.py --update from blueprint/audit.toml and the Lean environment. Do not edit.

\begin{imported}[\code{ConLeche.Denotes}]
  \label{imp:ConLeche-Denotes}
  \lean{ConLeche.LevelParam, ConLeche.BVarIdx, ConLeche.Denotes, ConLeche.Model}
  \leanok
  \srcloc{ConLeche/Denotes.lean}{101}
  \lead{In short} The set-model denotation relation (\code{Denotes}), with its valuations of bound
  variables (\code{BVarIdx}) and universe parameters (\code{LevelParam}).
  \lead{Used directly} 4 declarations: 2 definitions, 1 inductive type, 1 structure.
  \usedbystatements{\ref{def:EraseProof-sound_nonvacuous-statement},
    \ref{def:EraseProof-clCert_proof_pt-statement}}
\end{imported}

\begin{imported}[\code{ConLeche.Kernel.Core}]
  \label{imp:ConLeche-Kernel-Core}
  \lean{ConLeche.checkFuel}
  \leanok
  \srcloc{ConLeche/Kernel/Core.lean}{1925}
  \lead{In short} The fuel of the full-grade checker (\code{checkFuel}).
  \lead{Used directly} 1 declaration: 1 definition.
  \usedbystatements{\ref{def:EraseProof-clCert_proof_pt-statement}}
\end{imported}

\begin{imported}[\code{ConLeche.Kernel.Env}]
  \label{imp:ConLeche-Kernel-Env}
  \lean{ConLeche.CheckMode, ConLeche.Env}
  \leanok
  \srcloc{ConLeche/Kernel/Env.lean}{20}
  \lead{In short} Environments (\code{Env}) of declarations, and the checker modes
  (\code{CheckMode}).
  \lead{Used directly} 2 declarations: 1 inductive type, 1 structure.
  \usedbystatements{\ref{def:EraseProof-erase_correct-statement},
    \ref{def:EraseProof-erases_correct-statement}, \ref{def:EraseProof-srcEval_sound-statement},
    \ref{def:EraseProof-collectDeps_ok_of-statement},
    \ref{def:EraseProof-erase_correct_firstorder-statement},
    \ref{def:EraseProof-sound_nonvacuous-statement},
    \ref{def:EraseProof-clCert_proof_pt-statement}, \ref{def:EraseProof-collectDeps_sub-statement},
    \ref{def:EraseProof-markOf_agree-statement}}
\end{imported}

\begin{imported}[\code{ConLeche.Kernel.Error}]
  \label{imp:ConLeche-Kernel-Error}
  \lean{ConLeche.CheckError, ConLeche.CheckM}
  \leanok
  \srcloc{ConLeche/Kernel/Error.lean}{15}
  \lead{In short} The checker monad (\code{CheckM}) and its errors (\code{CheckError}).
  \lead{Used directly} 2 declarations: 1 definition, 1 inductive type.
  \usedbystatements{\ref{def:EraseProof-sound_nonvacuous-statement},
    \ref{def:EraseProof-clCert_proof_pt-statement}}
\end{imported}

\begin{imported}[\code{ConLeche.Kernel.Expr}]
  \label{imp:ConLeche-Kernel-Expr}
  \lean{ConLeche.Expr}
  \leanok
  \srcloc{ConLeche/Kernel/Expr.lean}{376}
  \lead{In short} The checker expressions (\code{Expr}): no binder names or infos, no
  metavariables.
  \lead{Used directly} 1 declaration: 1 inductive type.
  \usedbystatements{\ref{def:EraseProof-sound_nonvacuous-statement},
    \ref{def:EraseProof-clCert_proof_pt-statement}}
\end{imported}

\begin{imported}[\code{ConLeche.Kernel.TypeChecker}]
  \label{imp:ConLeche-Kernel-TypeChecker}
  \lean{ConLeche.whnf, ConLeche.inferTypeCore}
  \leanok
  \srcloc{ConLeche/Kernel/TypeChecker.lean}{33}
  \lead{In short} The weak-head reduction (\code{whnf}) and type inference (\code{inferTypeCore}).
  \lead{Used directly} 2 declarations: 2 definitions.
  \usedbystatements{\ref{def:EraseProof-sound_nonvacuous-statement},
    \ref{def:EraseProof-clCert_proof_pt-statement}}
\end{imported}

\begin{imported}[\code{ConLeche.Model.Denotes}]
  \label{imp:ConLeche-Model-Denotes}
  \lean{ConLeche.Model.Model.ofEnvModelM}
  \leanok
  \srcloc{ConLeche/Model/Denotes.lean}{328}
  \lead{In short} The model of an environment as a model of expressions
  (\code{Model.ofEnvModelM}).
  \lead{Used directly} 1 declaration: 1 definition.
  \usedbystatements{\ref{def:EraseProof-sound_nonvacuous-statement},
    \ref{def:EraseProof-clCert_proof_pt-statement}}
\end{imported}

\begin{imported}[\code{ConLeche.Model.Env.EnvModelM}]
  \label{imp:ConLeche-Model-Env-EnvModelM}
  \lean{ConLeche.Model.EnvModelM}
  \leanok
  \srcloc{ConLeche/Model/Env/EnvModelM.lean}{242}
  \lead{In short} Models of an accepted environment (\code{Model.EnvModelM}), the invariant behind
  the soundness theorem.
  \lead{Used directly} 1 declaration: 1 structure.
  \usedbystatements{\ref{def:EraseProof-sound_nonvacuous-statement},
    \ref{def:EraseProof-clCert_proof_pt-statement}}
\end{imported}

\begin{imported}[\code{ConLeche.SetTheory.Core}]
  \label{imp:ConLeche-SetTheory-Core}
  \lean{ConLeche.SetTheory}
  \leanok
  \srcloc{ConLeche/SetTheory/Core.lean}{95}
  \lead{In short} The class of set-theoretic universes (\code{SetTheory}), which the meaning claim
  quantifies over.
  \lead{Used directly} 1 declaration: 1 structure.
  \usedbystatements{\ref{def:EraseProof-sound_nonvacuous-statement},
    \ref{def:EraseProof-clCert_proof_pt-statement}}
\end{imported}

\begin{imported}[\code{ConLeche.SetTheory.Derive.Pt}]
  \label{imp:ConLeche-SetTheory-Derive-Pt}
  \lean{ConLeche.SetTheory.pt}
  \leanok
  \srcloc{ConLeche/SetTheory/Derive/Pt.lean}{68}
  \lead{In short} The proof point of the model (\code{SetTheory.pt}).
  \lead{Used directly} 1 declaration: 1 definition.
  \usedbystatements{\ref{def:EraseProof-sound_nonvacuous-statement},
    \ref{def:EraseProof-clCert_proof_pt-statement}}
\end{imported}
